\section{Jacobi Fields, Hessians, and Riccati Equations}

Let $(M,g)$ be an $n$-dimensional Riemannian manifold, let $U$ be a normal
neighborhood of $p\in M$, and let $r$ be the radial distance on $U$.  On
$U\setminus\{p\}$, the Gauss lemma gives $\operatorname{grad}r=\partial_r$.
Define the covariant Hessian $\nabla^2r$ and its raised-index operator
$\mathcal H_r=(\nabla^2r)^\sharp$ by
\begin{equation}
\mathcal H_r(w)=\nabla_w\partial_r. \tag{11.1}
\end{equation}
Thus
\begin{equation}
\langle\mathcal H_r(v),w\rangle_g=(\nabla^2r)(v,w). \tag{11.2}
\end{equation}
The endomorphism $\mathcal H_r$ is the \emph{Hessian operator} of $r$.

\begin{lemma}[Algebraic Properties of the Hessian Operator]
\label{lem:hessian-operator-properties}
\uses{prop:jacobi-field-hessian-equation, lem:symmetric-bilinear-form-unit-vectors}
\dcref{ch11:11.1}

The operator $\mathcal H_r$ is self-adjoint, satisfies
$\mathcal H_r(\partial_r)=0$, and restricts on each level set of $r$ to the
shape operator with normal $N=-\partial_r$.
\end{lemma}

\begin{proof}
\sketch
Symmetry of $\nabla^2r$ and (11.2) give self-adjointness.  Since radial integral
curves are geodesics, $\nabla_{\partial_r}\partial_r=0$.  Finally,
$\partial_r$ is a unit normal to each level set, so the Weingarten equation
identifies $\nabla\partial_r$ on its tangent space with the shape operator for
$-\partial_r$.
\end{proof}

\begin{proposition}[Jacobi Fields and the Hessian Operator]
\label{prop:jacobi-field-hessian-equation}
\uses{prop:jacobi-fields-vanishing-at-point, prop:normal-jacobi-equivalent-conditions}
\dcref{ch11:11.2}

Let $U$ be a normal neighborhood of $p$ in a Riemannian manifold, let
$r$ be its radial distance, and let $\gamma:[0,b]\to U$ be a unit-speed radial
geodesic.  If $J$ is a normal Jacobi field along $\gamma$ with $J(0)=0$, then
for $0<t\leq b$,
\begin{equation}
D_tJ(t)=\mathcal H_r(J(t)). \tag{11.3}
\end{equation}
\end{proposition}

\begin{proof}
\sketch
Write $v=\gamma'(0)$ and $w=D_tJ(0)\perp v$.  Choose a curve
$\sigma(s)\subset T_pM$ on the unit sphere with $\sigma(0)=v$ and
$\sigma'(0)=w$, and set
\[
\Gamma(s,t)=\exp_p(t\sigma(s)).
\]
Every main curve is a unit-speed radial geodesic, so
\begin{equation}
\partial_t\Gamma(s,t)=\partial_r|_{\Gamma(s,t)}. \tag{11.4}
\end{equation}
The variation field is $J$, and the symmetry lemma gives
\begin{equation}
D_tJ(t)=D_t\partial_s\Gamma(0,t)=D_s\partial_t\Gamma(0,t). \tag{11.5}
\end{equation}
The final term is $\nabla_{J(t)}\partial_r=\mathcal H_r(J(t))$.
\end{proof}

\begin{proposition}[Hessian of Radial Distance in Constant Curvature]
\label{prop:hessian-radial-distance-constant-curvature}
\dcref{ch11:11.3}
Let $U$ be a normal neighborhood of $p$ in a Riemannian manifold.  Its metric
has constant sectional curvature $c$ if and only if, on $U\setminus\{p\}$,
\begin{equation}
\mathcal H_r=\frac{s_c'(r)}{s_c(r)}\,\pi_r, \tag{11.6}
\end{equation}
where $s_c$ is defined by (10.8), and $\pi_r$ is orthogonal projection onto
the tangent space of the $r$-level set.
\end{proposition}

\begin{proof}
\sketch
Assume constant curvature.  Along a radial geodesic to $q$ with $r(q)=b$,
choose a parallel orthonormal frame with $E_n=\partial_r$.  The normal Jacobi
fields $J_i=s_cE_i$ vanish at $0$; absence of conjugate points in a normal
neighborhood implies $s_c(t)\ne0$ on $(0,b]$ (for $c=1/R^2>0$, this means
$b<\pi R$).  Applying (11.3) gives
\[
\mathcal H_r(E_i(b))=\frac{s_c'(b)}{s_c(b)}E_i(b),
\qquad
\mathcal H_r(E_n(b))=0,
\]
which is (11.6).

Conversely, (11.3) and (11.6) imply
$D_tJ=(s_c'/s_c)J$ for every normal Jacobi field vanishing at $0$; hence
$s_c^{-1}J$ is parallel and $J=ks_cE$.  The same normal-coordinate argument
as in Theorem 10.14 then gives the constant-curvature metric formula.
\end{proof}

For reference,
\[
\frac{s_c'(t)}{s_c(t)}=
\begin{cases}
1/t,&c=0,\\
(1/R)\cot(t/R),&c=1/R^2>0,\\
(1/R)\coth(t/R),&c=-1/R^2<0.
\end{cases}
\]

\begin{theorem}[The Riccati Equation]
\label{thm:riccati-equation}
\dcref{ch11:11.4}
Let $\gamma:[0,b]\to U$ be a unit-speed radial geodesic in a normal
neighborhood with radial Hessian operator $\mathcal H_r$.  Along
$\gamma|_{(0,b]}$,
\begin{equation}
D_t\mathcal H_r+\mathcal H_r^2+R_{\gamma'}=0, \tag{11.7}
\end{equation}
where $\mathcal H_r^2(w)=\mathcal H_r(\mathcal H_r(w))$ and
$R_{\gamma'}(w)=R(w,\gamma')\gamma'$.
\end{theorem}

\begin{theorem}[Riccati Comparison Theorem]
\label{thm:riccati-comparison-theorem}
\dcref{ch11:11.5}
Let $\gamma:[a,b]\to M$ be unit-speed.  Suppose self-adjoint endomorphism
fields $\eta,\widetilde\eta$ satisfy
\begin{equation}
D_t\eta+\eta^2+\sigma=0,
\qquad
D_t\widetilde\eta+\widetilde\eta^2+\widetilde\sigma=0, \tag{11.8}
\end{equation}
where $\sigma,\widetilde\sigma$ are continuous self-adjoint fields with
\begin{equation}
\widetilde\sigma(t)\geq\sigma(t)\quad(a\leq t\leq b). \tag{11.9}
\end{equation}
If $\lim_{t\searrow a}(\widetilde\eta-\eta)$ exists and is nonpositive, then
\[
\widetilde\eta(t)\leq\eta(t)\qquad(a<t\leq b).
\]
\end{theorem}

In a parallel orthonormal frame, compare self-adjoint endomorphisms by their
quadratic forms: $A\leq B$ means
$\langle Av,v\rangle\leq\langle Bv,v\rangle$ for all $v$.  Every square of a
self-adjoint endomorphism is positive semidefinite.

\begin{theorem}[Matrix Riccati Comparison Theorem]
\label{thm:matrix-riccati-comparison}
\dcref{ch11:11.6}
Let $H,\widetilde H:[a,b]\to S(n,\mathbb R)$ satisfy
\begin{equation}
H'+H^2+S=0,
\qquad
\widetilde H'+\widetilde H^2+\widetilde S=0, \tag{11.10}
\end{equation}
where $S,\widetilde S$ are continuous symmetric matrices and
\begin{equation}
\widetilde S(t)\geq S(t). \tag{11.11}
\end{equation}
If
\begin{equation}
\lim_{t\searrow a}(\widetilde H(t)-H(t))\leq0, \tag{11.12}
\end{equation}
then
\begin{equation}
\widetilde H(t)\leq H(t)\qquad(a<t\leq b). \tag{11.13}
\end{equation}
\end{theorem}

\begin{proof}
\sketch
Set
\[
A=\widetilde H-H,
\qquad
B=-\tfrac12(\widetilde H+H).
\]
Then $A$ extends continuously to $a$ with $A(a)\leq0$, and
\begin{align*}
A'&=BA+AB-\widetilde S+S,\\
B'&=\tfrac12(\widetilde H^2+\widetilde S+H^2+S).
\end{align*}
Thus
\begin{equation}
A'\leq BA+AB, \tag{11.14}
\end{equation}
and, for some constant $k$,
\begin{equation}
B'\geq k\operatorname{Id}. \tag{11.15}
\end{equation}
Integrating gives
\begin{equation}
B(t)=B(b)-\int_t^bB'(u)\,du, \tag{11.16}
\end{equation}
so $B$ is bounded above.  Choose $K$ with $B-K\operatorname{Id}$ negative
definite and define
\[
f(t,x)=e^{-2Kt}\langle A(t)x,x\rangle,
\qquad x\in S^{n-1}.
\]
If $f$ had a positive maximum $(t_0,x_0)$, then $x_0$ would be an eigenvector
of $A(t_0)$ with positive eigenvalue $\lambda_0$, and the left derivative at
$t_0$ would be nonnegative.  But (11.14) gives
\[
\partial_tf(t_0,x_0)
\leq2e^{-2Kt_0}\lambda_0
\langle x_0,(B(t_0)-K\operatorname{Id})x_0\rangle<0,
\]
a contradiction.  Hence $A\leq0$, proving (11.13).
\end{proof}

\begin{proof}
\sketch
Choose a parallel orthonormal frame $(E_i)$ along $\gamma$ and represent
$\eta,\widetilde\eta,\sigma,\widetilde\sigma$ by symmetric matrices
$H,\widetilde H,S,\widetilde S$.  Since $D_tE_i=0$, (11.8) becomes (11.10).
The matrix comparison theorem gives $\widetilde H\leq H$, hence
$\widetilde\eta\leq\eta$.
\end{proof}
\section{Comparisons Based on Sectional Curvature}

\begin{theorem}[Hessian Comparison]
\label{thm:hessian-comparison}
\uses{prop:hessian-radial-distance-constant-curvature}
\dcref{ch11:11.7}

Let $(M,g)$ be a Riemannian $n$-manifold, let $p\in M$, let $U$ be a normal
neighborhood of $p$, and let $r$ be its radial distance function.
\begin{enumerate}
\item[(a)] If $\sec_M\leq c$, then on $U_0\setminus\{p\}$,
\begin{equation}
\mathcal H_r\geq\frac{s_c'(r)}{s_c(r)}\pi_r, \tag{11.17}
\end{equation}
where $U_0=U$ for $c\leq0$, while for $c=1/R^2>0$,
\[
U_0=\{q\in U:r(q)<\pi R\}.
\]
\item[(b)] If $\sec_M\geq c$, then on all of $U\setminus\{p\}$,
\begin{equation}
\mathcal H_r\leq\frac{s_c'(r)}{s_c(r)}\pi_r. \tag{11.18}
\end{equation}
\end{enumerate}
\end{theorem}

\begin{proof}
\sketch Let $U_0$ be the set where $s_c(r)>0$ and define
\[
\mathcal H_r^c=\frac{s_c'(r)}{s_c(r)}\pi_r.
\]
Along a unit-speed radial geodesic $\gamma\colon[0,b]\to U_0$, the orthogonal
projection
\[
\pi_r(w)=w-\langle w,\gamma'\rangle\gamma'
\]
extends smoothly to $t=0$ and satisfies $D_t\pi_r=0$. Since
$s_c''=-cs_c$ and $\pi_r^2=\pi_r$,
\[
D_t\mathcal H_r^c+(\mathcal H_r^c)^2+c\pi_r=0.
\]
The radial Hessian satisfies
\[
D_t\mathcal H_r+\mathcal H_r^2+R_{\gamma'}=0.
\]
The curvature hypotheses give
\[
R_{\gamma'}\leq c\pi_r\quad\text{when }\sec_M\leq c,
\qquad
R_{\gamma'}\geq c\pi_r\quad\text{when }\sec_M\geq c.
\]

The initial singularities cancel. Namely,
\begin{equation}
\frac{s_c'(r)}{s_c(r)}=\frac1r+O(r), \tag{11.19}
\end{equation}
and in normal coordinates,
\begin{equation}
\mathcal H_r=\frac1r\pi_r+O(r). \tag{11.20}
\end{equation}
For (11.20), use
\[
\mathcal H_r
=g^{ij}\bigl(\partial_j\partial_kr
-\Gamma_{jk}^m\partial_mr\bigr)\partial_i\otimes dx^k,
\]
together with
\[
\partial_mr=\frac{x^m}{r},
\quad g^{ij}=\delta^{ij}+O(r^2),
\quad \Gamma_{jk}^m=O(r).
\]
The Euclidean leading term is $r^{-1}\pi_r$. Thus
$\mathcal H_r-\mathcal H_r^c\to0$ as $t\searrow0$, and Riccati comparison gives
(11.17) and (11.18) on $U_0$.

For the lower-curvature case, $U_0=U$. This is automatic when $c\leq0$. If
$c=1/R^2>0$, then $s_c'/s_c\to-\infty$ as $r\nearrow\pi R$. If a radial
segment in the normal neighborhood reached $r=\pi R$, inequality (11.18)
would contradict smoothness of $\mathcal H_r$ there.
\end{proof}

\begin{corollary}[Principal Curvature Comparison]
\label{cor:principal-curvature-comparison}
\uses{thm:hessian-comparison, prop:hessian-radial-distance-constant-curvature, lem:hessian-operator-properties}
\dcref{ch11:11.8}

Under the notation of the preceding theorem, the principal curvatures of the
$r$-level sets with respect to the inward unit normal satisfy:
\begin{enumerate}
\item[(a)] if $\sec_M\leq c$, then on $U_0\setminus\{p\}$,
\[
\kappa\geq\frac{s_c'(r)}{s_c(r)},
\]
with the same definition of $U_0$ as above;
\item[(b)] if $\sec_M\geq c$, then on $U\setminus\{p\}$,
\[
\kappa\leq\frac{s_c'(r)}{s_c(r)}.
\]
\end{enumerate}

\begin{proof}
\sketch The shape operator of an $r$-level set for the inward normal is the
restriction of $\mathcal H_r$ to its tangent space. Restrict the operator
inequalities in the Hessian comparison theorem and evaluate on eigenvectors.
\end{proof}
\end{corollary}

\begin{theorem}[Jacobi Field Comparison]
\label{thm:jacobi-field-comparison}
\uses{cor:jacobi-field-decomposition}
\dcref{ch11:11.9}

Let $\gamma\colon[0,b]\to M$ be a unit-speed geodesic and let $J$ be a normal
Jacobi field with $J(0)=0$.
\begin{enumerate}
\item[(a)] If $\sec_M\leq c$, then
\begin{equation}
|J(t)|\geq s_c(t)|D_tJ(0)| \tag{11.21}
\end{equation}
for $0\leq t\leq b_1$, where $b_1=b$ for $c\leq0$ and
$b_1=\min(b,\pi R)$ for $c=1/R^2>0$.
\item[(b)] If $\sec_M\geq c$, then
\begin{equation}
|J(t)|\leq s_c(t)|D_tJ(0)| \tag{11.22}
\end{equation}
for $0\leq t\leq b_2$, where $\gamma(b_2)$ is the first conjugate point to
$\gamma(0)$ along the segment if one exists, and $b_2=b$ otherwise.
\end{enumerate}
\end{theorem}

\begin{proof}
\sketch The result is immediate if $D_tJ(0)=0$, so assume otherwise. Let
$b_0$ be the largest time such that there is no conjugate point in $(0,b_0)$
and $s_c(t)>0$ there. First suppose $\gamma([0,b_0))$ lies in a normal
neighborhood of $p=\gamma(0)$. On $(0,b_0)$, define
\[
f(t)=\log\frac{|J(t)|}{s_c(t)}.
\]
Since $D_tJ=\mathcal H_r(J)$,
\[
f'(t)
=\frac{\langle\mathcal H_rJ,J\rangle}{|J|^2}
-\frac{s_c'}{s_c}.
\]
Hessian comparison gives $f'\geq0$ under an upper curvature bound and
$f'\leq0$ under a lower curvature bound. Moreover, two applications of
l'Hopital's rule and the Jacobi equations give
\begin{align*}
\lim_{t\searrow0}\frac{|J|^2}{s_c^2}
&=\lim_{t\searrow0}
\frac{-2\langle R(J,\gamma')\gamma',J\rangle+2|D_tJ|^2}
{2s_c''s_c+2(s_c')^2}\\
&=|D_tJ(0)|^2.
\end{align*}
Thus the corresponding inequality holds through $b_0$.

For a general segment, put $v=\gamma'(0)$. On a convex open neighborhood
$W\subseteq T_pM$ of the ray $\{tv:0\leq t<b_0\}$ on which $\exp_p$ is a
local diffeomorphism, use the pullback metric
$\widetilde g=\exp_p^*g$. Then $\exp_p$ is a local isometry,
$\widetilde\gamma(t)=tv$ is radial, and
\[
\widetilde J(t)=(d\exp_p)_{tv}^{-1}J(t)
\]
is a Jacobi field along $\widetilde\gamma$. The normal-neighborhood argument
therefore yields the same estimate for $J$ on $[0,b_0]$.

In case (a), if $b_0<b_1$, a nonzero Jacobi field vanishing at $0$ and $b_0$
would satisfy
\[
0=|J(b_0)|\geq s_c(b_0)|D_tJ(0)|>0,
\]
a contradiction. Hence $b_0\geq b_1$. In case (b), if $b_0<b_2$, then
$s_c(b_0)=0$ while $\gamma(b_0)$ is not conjugate to $\gamma(0)$. Any nonzero
normal Jacobi field vanishing at $0$ would satisfy $|J(b_0)|=0$, again a
contradiction. Thus $b_0\geq b_2$.
\end{proof}

\begin{theorem}[Metric Comparison]
\label{thm:metric-comparison}
\uses{thm:gauss-lemma, cor:jacobi-field-tangent-vectors, prop:jacobi-fields-vanishing-at-point, prop:jacobi-constant-curvature, thm:jacobi-field-comparison, cor:polar-decomposition-integrals}
\dcref{ch11:11.10}

Let $(U,(x^i))$ be a normal coordinate chart for $g$ centered at $p$, and let
$g_c$ be the constant-curvature metric in the same coordinates as in (10.17).
\begin{enumerate}
\item[(a)] If $\sec_g\leq c$, then
\[
g(w,w)\geq g_c(w,w)
\]
for every $q\in U\setminus\{p\}$ and $w\in T_qM$, with the additional
restriction $d_g(p,q)<\pi R$ when $c=1/R^2>0$.
\item[(b)] If $\sec_g\geq c$, then
\[
g(w,w)\leq g_c(w,w)
\]
for all $q\in U\setminus\{p\}$ and $w\in T_qM$.
\end{enumerate}
\end{theorem}

\begin{proof}
\sketch Put $b=d_g(p,q)$ and decompose
$w=y+z$ into radial and level-set components. The Gauss lemma for both metrics
gives
\[
g(w,w)=g(y,y)+g(z,z),
\qquad
g_c(w,w)=g_c(y,y)+g_c(z,z).
\]
The radial field has unit length for both metrics, so the radial terms agree.

Let $\gamma\colon[0,b]\to U$ be radial. There is a normal Jacobi field $J$
with $J(0)=0$ and $J(b)=z$, having normal-coordinate expression
\[
J(t)=t a^i\partial_i|_{\gamma(t)}.
\]
The same coordinate field is the corresponding Jacobi field for $g_c$, and
the two metrics agree at $p$. Under the upper bound,
\begin{align*}
g(z,z)=|J(b)|_g^2
&\geq s_c(b)^2|D_tJ(0)|_g^2\\
&=s_c(b)^2|D_tJ(0)|_{g_c}^2
=|J(b)|_{g_c}^2=g_c(z,z).
\end{align*}
The lower-bound case uses the reversed Jacobi comparison inequality.
\end{proof}

\begin{theorem}[Laplacian Comparison 1]
\label{thm:laplacian-comparison-1}
\uses{thm:hessian-comparison, prop:hessian-radial-distance-constant-curvature}
\dcref{ch11:11.11}

Let $(M,g)$ be a Riemannian $n$-manifold with $\sec_M\leq c$, let $U$ be a
normal neighborhood of $p$, and let $r$ be radial distance. On
$U_0\setminus\{p\}$,
\[
\Delta r\geq(n-1)\frac{s_c'(r)}{s_c(r)},
\]
where $U_0=U$ for $c\leq0$ and
$U_0=\{q\in U:r(q)<\pi R\}$ for $c=1/R^2>0$.
\end{theorem}

\begin{proof}
\sketch Since
\[
\Delta r=\operatorname{tr}_g(\nabla^2r)=\operatorname{tr}(\mathcal H_r),
\qquad \operatorname{tr}(\pi_r)=n-1,
\]
the result is the trace of the upper-curvature Hessian comparison inequality.
\end{proof}

\begin{theorem}[Conjugate Point Comparison I]
\label{thm:conjugate-point-comparison-I}
\lean{LeeLib.Ch11.conjugatePointComparisonI_nonpos, LeeLib.Ch11.conjugatePointComparisonI_pos}
\uses{thm:jacobi-field-comparison}
\leanok
\dcref{ch11:11.12}



Let $(M,g)$ be a Riemannian manifold with $\sec_M\leq c$. If $c\leq0$, no
point has a conjugate point along any geodesic. If $c=1/R^2>0$, no geodesic
segment of length less than $\pi R$ contains conjugate endpoints.

\begin{proof}
\sketch Let $J$ be a nonzero normal Jacobi field along a unit-speed geodesic,
with $J(0)=0$. Then $D_tJ(0)\ne0$. If $c\leq0$, comparison with $c=0$ gives
\[
|J(t)|\geq t|D_tJ(0)|>0
\]
for $t>0$. If $c=1/R^2>0$, Jacobi comparison gives
\[
|J(t)|\geq R\sin(t/R)|D_tJ(0)|>0
\]
for $0<t<\pi R$. Thus $J$ cannot vanish again in the stated ranges.
\end{proof}
\end{theorem}

\begin{lemma}
\label{lem:riemannian-volume-laplacian}
\uses{cor:gradient-radial-distance, prop:coordinate-representations-divergence-laplacian}
\dcref{ch11:11.13}

Let $(x^i)$ be Riemannian normal coordinates on a normal neighborhood $U$ of
$p$, let $r$ be radial distance, and let $\det g=\det(g_{ij})$. On
$U\setminus\{p\}$,
\begin{equation}
\Delta r=\partial_r\log\bigl(r^{n-1}\sqrt{\det g}\bigr). \tag{11.23}
\end{equation}

\begin{proof}
\sketch The Gauss lemma gives
\[
\operatorname{grad}r=\partial_r,
\qquad
g^{ij}\partial_jr=\frac{x^i}{r}.
\]
The coordinate Laplacian formula yields
\begin{align*}
\Delta r
&=\frac1{\sqrt{\det g}}
  \partial_i\left(g^{ij}\sqrt{\det g}\,\partial_jr\right)\\
&=\partial_i\left(\frac{x^i}{r}\right)
  +\frac{x^i}{r\sqrt{\det g}}\partial_i\sqrt{\det g}\\
&=\frac{n-1}{r}
  +\frac1{\sqrt{\det g}}\partial_r\sqrt{\det g},
\end{align*}
which is (11.23).
\end{proof}
\end{lemma}

G\"unther's comparison theorem \cite{Gun60} gives the sectional-curvature
volume estimate and its equality case.

\begin{theorem}[G\"unther's Volume Comparison]
\label{thm:gunther-volume-comparison}
\uses{thm:laplacian-comparison-1, lem:riemannian-volume-laplacian, cor:polar-decomposition-integrals, thm:hessian-comparison, prop:hessian-radial-distance-constant-curvature, thm:bishop-gromov-volume-comparison}
\dcref{ch11:11.14}

Let $(M,g)$ be connected with $\sec_M\leq c$, fix $p\in M$, and set
\[
\delta_0=
\begin{cases}
\operatorname{inj}(p),&c\leq0,\\
\min(\pi R,\operatorname{inj}(p)),&c=1/R^2>0.
\end{cases}
\]
Let $V_g(\delta)$ be the volume of $B_\delta(p)$, and let $V_c(\delta)$ be the
volume of the radius-$\delta$ ball in the simply connected $n$-dimensional
constant-curvature-$c$ model. For $0<\delta\leq\delta_0$,
\[
V_g(\delta)\geq V_c(\delta). \qquad (11.24)
\]
The ratio $V_g(\delta)/V_c(\delta)$ is nondecreasing and tends to $1$ as
$\delta\searrow0$. If equality holds for some
$\delta\in(0,\delta_0]$, then $g$ has constant sectional curvature $c$ on
$B_\delta(p)$.

\begin{proof}
\sketch Use normal coordinates on $B_{\delta_0}(p)$ and let $g_c$ be the
constant-curvature metric in the same coordinates. Laplacian comparison and
the preceding lemma give
\[
\partial_r\log\bigl(r^{n-1}\sqrt{\det g}\bigr)
=\Delta r
\geq(n-1)\frac{s_c'(r)}{s_c(r)}
=\partial_r\log\bigl(s_c(r)^{n-1}\bigr).
\qquad (11.25)
\]
Hence along each radial direction,
\[
\lambda(r,\omega)
=\frac{r^{n-1}\sqrt{\det g(r,\omega)}}{s_c(r)^{n-1}}
\]
is nondecreasing. Since $g_{ij}(0)=\delta_{ij}$ and $s_c(r)/r\to1$, the
convergence $\lambda(r,\omega)\to1$ as $r\searrow0$ is uniform in $\omega$.

Polar integration gives
\[
V_g(\delta)
=\int_{S^{n-1}}\int_0^\delta
\lambda(\rho,\omega)s_c(\rho)^{n-1}\,d\rho\,d\omega
\]
and
\[
V_c(\delta)
=\operatorname{Vol}(S^{n-1})
  \int_0^\delta s_c(\rho)^{n-1}\,d\rho.
\]
Thus
\[
\frac{V_g(\delta)}{V_c(\delta)}
=\frac1{\operatorname{Vol}(S^{n-1})}
\int_{S^{n-1}}
\frac{\int_0^\delta
\lambda(\rho,\omega)s_c(\rho)^{n-1}\,d\rho}
{\int_0^\delta s_c(\rho)^{n-1}\,d\rho}
\,d\omega.
\qquad (11.26)
\]
For each $\omega$, the quotient inside the integral is the weighted average of
the nondecreasing function $\lambda(\cdot,\omega)$ over $[0,\delta]$; splitting
the integrals at $0<\delta_1<\delta_2$ shows that this average is
nondecreasing in $\delta$. Its limit is $1$, proving monotonicity and
$V_g\geq V_c$.

If $V_g(\delta)=V_c(\delta)$, then monotonicity and continuity force
$\lambda\equiv1$ on $(0,\delta)\times S^{n-1}$. Equality holds in (11.25), so
\[
\operatorname{tr}(\mathcal H_r)
=\operatorname{tr}\left(\frac{s_c'(r)}{s_c(r)}\pi_r\right).
\]
Hessian comparison makes
$\mathcal H_r-(s_c'/s_c)\pi_r$ positive semidefinite. Its trace is zero, so
the endomorphism vanishes. The constant-model radial Hessian characterization
then gives sectional curvature identically $c$ on $B_\delta(p)$.
\end{proof}
\end{theorem}
\section{Comparisons Based on Ricci Curvature}

\begin{theorem}[Laplacian Comparison II]
\label{thm:laplacian-comparison-ii}
\uses{thm:laplacian-comparison-1, cor:jacobi-field-decomposition}
\dcref{ch11:11.15}


Let $(M,g)$ be a Riemannian $n$-manifold satisfying
\[
\operatorname{Rc}(v,v)\geq(n-1)c
\]
for every unit vector $v$. For $p\in M$, let $U$ be a normal neighborhood,
and let $r$ be its radial distance function. Then on
$U_0\setminus\{p\}$,
\begin{equation}
\Delta r\leq(n-1)\frac{s_c'(r)}{s_c(r)},
\tag{11.27}
\end{equation}
where $s_c$ is defined in (10.8), and
\[
U_0=
\begin{cases}
U,&c\leq0,\\
\{q\in U:r(q)<\pi R\},&c=1/R^2>0.
\end{cases}
\]
\end{theorem}

\begin{proof}
\sketch Fix a unit-speed radial geodesic
$\gamma:[0,b]\to U_0$. Tracing the Riccati equation for the radial Hessian
$\mathcal H_r$ gives
\begin{equation}
\frac d{dt}\Delta r
+\operatorname{tr}(\mathcal H_r^2)
+\operatorname{tr}(R_{\gamma'})=0.
\tag{11.28}
\end{equation}
In an orthonormal frame,
$\operatorname{tr}(R_{\gamma'})=\operatorname{Rc}(\gamma',\gamma')$.
Define the trace-free radial Hessian
\[
\widehat{\mathcal H}_r
=\mathcal H_r-\frac{\Delta r}{n-1}\pi_r,
\]
where $\pi_r$ projects orthogonally onto $\partial_r^\perp$. Since
$\mathcal H_r\partial_r=0$, $\mathcal H_r$ is self-adjoint, and
$\pi_r^2=\pi_r$,
\[
\operatorname{tr}(\widehat{\mathcal H}_r^2)
=\operatorname{tr}(\mathcal H_r^2)-\frac{(\Delta r)^2}{n-1}.
\]
Thus
\begin{equation}
\frac d{dt}\left(\frac{\Delta r}{n-1}\right)
+\left(\frac{\Delta r}{n-1}\right)^2
+\frac{
\operatorname{tr}(\widehat{\mathcal H}_r^2)
+\operatorname{Rc}(\gamma',\gamma')
}{n-1}=0.
\tag{11.30}
\end{equation}

Set
\[
H(t)=\frac{s_c'(t)}{s_c(t)},\qquad
\widetilde H(t)=\frac{\Delta r}{n-1}\bigg|_{\gamma(t)}.
\]
The model function satisfies $H'+H^2+c=0$. Since
$\widehat{\mathcal H}_r^2$ is positive semidefinite and the Ricci lower bound
holds,
\[
\widetilde S(t)
=\frac{
\operatorname{tr}(\widehat{\mathcal H}_r^2)
+\operatorname{Rc}(\gamma',\gamma')
}{n-1}
\geq c.
\]
The expansions
\[
\frac{s_c'(r)}{s_c(r)}=\frac1r+O(r),
\qquad
\mathcal H_r=\frac1r\pi_r+O(r)
\]
imply
$\widehat{\mathcal H}_r\to0$ and
$\widetilde H-H\to0$ as $t\searrow0$, while
$\widetilde S$ extends continuously with limiting value at least $c$.
Scalar Riccati comparison yields $\widetilde H\leq H$ on $(0,b]$, proving
(11.27).
\end{proof}

\begin{theorem}[Conjugate Point Comparison II]
\label{thm:conjugate-points-ricci-bounded}
\uses{thm:laplacian-comparison-ii, lem:riemannian-volume-laplacian, thm:jacobi-field-comparison}
\dcref{ch11:11.16}


Let $(M,g)$ be a Riemannian $n$-manifold satisfying
\[
\operatorname{Rc}(v,v)\geq(n-1)c
\qquad\text{for all unit }v,
\]
where $c=1/R^2>0$. Every geodesic segment of length at least $\pi R$
contains a conjugate point.

\begin{proof}
\sketch Let $U$ be a normal neighborhood of $p$. On
$U_0=\{r<\pi R\}$, Laplacian comparison and the radial volume-density formula
give
\[
\partial_r\log\bigl(r^{n-1}\sqrt{\det g}\bigr)
=\Delta r
\leq(n-1)\frac{s_c'(r)}{s_c(r)}
=\partial_r\log(s_c(r)^{n-1}).
\]
Because
$r^{n-1}\sqrt{\det g}/s_c(r)^{n-1}\to1$ as $r\to0$,
\begin{equation}
\sqrt{\det g}\leq\frac{s_c(r)^{n-1}}{r^{n-1}}.
\tag{11.31}
\end{equation}
Since $s_c(\pi R)=0$, a normal neighborhood cannot contain a radial point
with $r\geq\pi R$, because (11.31) would force its normal-coordinate metric
determinant to vanish.

Now let $\gamma:[0,b]\to M$ be unit speed with $b\geq\pi R$. If it had no
conjugate point, $\exp_{\gamma(0)}$ would be a local diffeomorphism on a
star-shaped neighborhood of the radial segment in the tangent space. Pull
back $g$ there. The pulled-back radial geodesic has length at least $\pi R$,
the same Ricci lower bound, and lies in a normal neighborhood, contradicting
the preceding paragraph.
\end{proof}
\end{theorem}

\begin{corollary}[Injectivity Radius Comparison]
\label{cor:injectivity-radius-ricci-bounded}
\uses{thm:conjugate-points-ricci-bounded}
\dcref{ch11:11.17}


Under the hypotheses of Theorem~11.16,
\[
\operatorname{inj}(p)\leq\pi R
\qquad(p\in M).
\]

\begin{proof}
\sketch Radial geodesics in a geodesic ball are minimizing, while a geodesic
segment containing a conjugate point cannot minimize beyond it. Theorem~11.16
therefore bounds every geodesic-ball radius by $\pi R$.
\end{proof}
\end{corollary}

\begin{corollary}[Diameter Comparison]
\label{cor:diameter-ricci-bounded}
\uses{thm:conjugate-points-ricci-bounded}
\dcref{ch11:11.18}


Let $(M,g)$ be a complete, connected Riemannian $n$-manifold satisfying
$\operatorname{Rc}(v,v)\geq(n-1)/R^2$ for every unit vector. Then
\[
\operatorname{diam}(M)\leq\pi R.
\]

\begin{proof}
\sketch Completeness and connectedness provide a minimizing geodesic between
any two points. Theorem~11.16 bounds its length by $\pi R$.
\end{proof}
\end{corollary}

\begin{theorem}[Bishop--Gromov Volume Comparison]
\label{thm:bishop-gromov-volume-comparison}
\uses{thm:chengs-maximal-diameter, thm:gunther-volume-comparison, thm:laplacian-comparison-ii, thm:cut-locus-properties, cor:polar-decomposition-integrals}
\dcref{ch11:11.19}


Let $(M,g)$ be a connected Riemannian $n$-manifold with
\[
\operatorname{Rc}(v,v)\geq(n-1)c
\]
for every unit vector. Fix $p\in M$. Let $V_g(\delta)$ be the volume of the
metric $\delta$-ball about $p$, and let $V_c(\delta)$ be the volume of a
metric $\delta$-ball in the simply connected $n$-dimensional model of
constant sectional curvature $c$. For
$0<\delta\leq\operatorname{inj}(p)$,
\begin{equation}
V_g(\delta)\leq V_c(\delta).
\tag{11.32}
\end{equation}
Moreover, $V_g(\delta)/V_c(\delta)$ is nonincreasing in $\delta$ and tends to
$1$ as $\delta\searrow0$. If $(M,g)$ is complete, these conclusions hold for
all $\delta>0$. In either setting, equality in (11.32) at some $\delta$
implies that $g$ has constant sectional curvature $c$ throughout the metric
$\delta$-ball about $p$.

\begin{proof}
\sketch First suppose
$0<\delta\leq\operatorname{inj}(p)$, so the metric ball is geodesic.
Integrating Theorem~11.15 along each radial geodesic shows that the radial
volume-density ratio to the model is nonincreasing and tends to $1$ at the
center. Polar integration then makes
$V_g(\delta)/V_c(\delta)$ nonincreasing with limit $1$, proving (11.32).

Assume now that $M$ is complete. The cut locus has measure zero, and
$\exp_p$ maps $\operatorname{ID}(p)$ diffeomorphically onto its complement.
Thus
\[
V_g(\delta)
=\int_{\operatorname{ID}(p)\cap B_\delta(0)}
\sqrt{\det g}\,dV_{\widetilde g}.
\]
For $\Phi(\rho,\omega)=\rho\omega$, define
\[
\widetilde s_c(\rho)=
\begin{cases}
s_c(\rho),&c\leq0,\\
s_c(\rho),&c=1/R^2>0\text{ and }\rho<\pi R,\\
0,&c=1/R^2>0\text{ and }\rho\geq\pi R,
\end{cases}
\]
and
\[
\widetilde\lambda(\rho,\omega)=
\begin{cases}
\dfrac{
(\sqrt{\det g})\circ\Phi(\rho,\omega)\,\rho^{n-1}
}{
\widetilde s_c(\rho)^{n-1}
},
&\Phi(\rho,\omega)\in\operatorname{ID}(p),\\
0,&\Phi(\rho,\omega)\notin\operatorname{ID}(p).
\end{cases}
\]
In the positive-curvature case, diameter comparison ensures
$\widetilde s_c(\rho)>0$ on $\operatorname{ID}(p)$. Polar integration gives
\begin{align*}
V_g(\delta)
&=\int_{S^{n-1}}\int_0^\delta
\widetilde\lambda(\rho,\omega)
\widetilde s_c(\rho)^{n-1}\,d\rho\,dV_{\overline g},\\
V_c(\delta)
&=\int_{S^{n-1}}\int_0^\delta
\widetilde s_c(\rho)^{n-1}\,d\rho\,dV_{\overline g}.
\end{align*}
For every $\omega$, the Laplacian comparison argument makes
$\widetilde\lambda(\rho,\omega)$ nonincreasing in $\rho$. The same integral
comparison proves the complete-case monotonicity, limit, and volume bound.

Suppose equality holds at $\delta$ and first assume
$\delta\leq\operatorname{inj}(p)$. Equality forces the radial density ratio
to be identically $1$ for $0<r<\delta$, hence
\begin{equation}
\Delta r=(n-1)\frac{s_c'(r)}{s_c(r)}
\tag{11.33}
\end{equation}
on $B_\delta(p)\setminus\{p\}$. Along each radial geodesic,
$u=(\Delta r)/(n-1)=s_c'/s_c$ satisfies $u'+u^2+c=0$. Comparison with
(11.30) yields
\[
\frac{
\operatorname{tr}(\widehat{\mathcal H}_r^2)
+\operatorname{Rc}(\gamma',\gamma')
}{n-1}=c.
\]
Both the trace term and
$\operatorname{Rc}(\gamma',\gamma')-(n-1)c$ are nonnegative, so
$\widehat{\mathcal H}_r=0$ and
\[
\mathcal H_r
=\frac{\Delta r}{n-1}\pi_r
=\frac{s_c'(r)}{s_c(r)}\pi_r.
\]
The radial Hessian characterization then gives constant sectional curvature
$c$ on $B_\delta(p)$.

In the complete case, equality similarly forces
$\widetilde\lambda=1$ where $0<\rho<\delta$ and
$\widetilde s_c(\rho)>0$. If $c\leq0$, or if
$c=1/R^2>0$ and $\delta\leq\pi R$, this makes the metric ball geodesic and
the preceding argument applies. If $c=1/R^2>0$ and $\delta>\pi R$, diameter
comparison makes the $\delta$-ball all of $M$. Monotonicity gives equality at
$\pi R$, so $g$ has curvature $c$ on the $\pi R$-ball and hence, by
continuity, on its closure $M$.
\end{proof}
\end{theorem}

\begin{corollary}
\label{cor:ricci-volume-comparison}
\uses{thm:bishop-gromov-volume-comparison}
\dcref{ch11:11.20}


Let $(M,g)$ be a compact Riemannian $n$-manifold satisfying
$\operatorname{Rc}(v,v)\geq(n-1)/R^2$ for every unit vector. Then
\[
\operatorname{Vol}(M)
\leq\operatorname{Vol}\bigl(S^n(R)\bigr).
\]
If equality holds, $(M,g)$ has constant sectional curvature $1/R^2$.
\end{corollary}

\begin{exercise}\label{prob:focal-points-hypersurface}\dcref{ch11:11-6}
Suppose $P$ is an embedded hypersurface in a Riemannian manifold $(M, g)$ and $N$ is a unit normal vector field along $P$. Suppose the principal curvatures of $P$ with respect to $-N$ are bounded below by a constant $\lambda$, and the sectional curvatures of $M$ are bounded above by $-\lambda^2$. Prove that $P$ has no focal points along any geodesic segment with initial velocity $N_p$ for $p \in P$.
\end{exercise}

\begin{exercise}\label{prob:focal-points-hypersurface-ii}\dcref{ch11:11-7}
\uses{prob:transverse-jacobi-comparison}
Suppose $P$ is an embedded hypersurface in a Riemannian manifold $(M, g)$ and $N$ is a unit normal vector field along $P$. Suppose the sectional curvatures of $M$ are bounded below by a constant $c$, and the principal curvatures of $P$ with respect to $-N$ are bounded above by a constant $\lambda$. Let $u$ be the solution to the initial value problem (11.35). Prove that if $b$ is a positive real number such that $u(b) = 0$, then $P$ has a focal point along every geodesic segment with initial velocity $N_p$ for some $p \in P$ and with length greater than or equal to $b$.
\end{exercise}

\begin{exercise}
\label{prob:hessian-polar-coordinates}
\dcref{ch11:11-1}
\uses{cor:semigeodesic-metric-christoffel,ex:polar-normal-coordinates}

Let $(M,g)$ be a Riemannian manifold, and let $U$ be a normal neighborhood of $p \in M$. Use Corollary 6.42 to show that in every choice of polar normal coordinates $(\theta^1,\ldots,\theta^{n-1},r)$ on a subset of $U \setminus \{p\}$ (see Example 6.45), the covariant Hessian of $r$ is given by
\[
\nabla^2 r = \frac{1}{2} \sum_{\alpha,\beta=1}^{n-1} (\partial_r g_{\alpha\beta}) \, d\theta^\alpha d\theta^\beta.
\]
\end{exercise}

\begin{exercise}\label{prob:injectivity-radius-comparison}\dcref{ch11:11-4}
\uses{prob:distance-cut-locus}
Let $(M, g)$ be a compact Riemannian manifold. Prove that if $R, L$ are positive numbers such that all sectional curvatures of $M$ are less than or equal to $1/R^2$ and all closed geodesics have lengths greater than or equal to $L$, then
\[
\text{inj}(M) \geq \min(\pi R, \tfrac{1}{2}L).
\]
[Hint: Assume not, and use the result of Problem 10-23(b).]
\end{exercise}

\begin{exercise}
\label{prob:jacobi-field-submanifold}
\dcref{ch11:11-2}
\uses{prop:jacobi-field-hessian-equation,prop:fermi-coordinates-radial-function,prob:transverse-jacobi-field}

Prove the following extension to Proposition 11.2: Suppose $P$ is an embedded submanifold of a Riemannian manifold $(M,g)$, $U$ is a normal neighborhood of $P$ in $M$, and $r$ is the radial distance function for $P$ in $U$ (see Prop.\ 6.37). If $\gamma : [0,b] \to U$ is a geodesic segment with $\gamma(0) \in P$ and $\gamma'(0)$ normal to $P$, and $J$ is a Jacobi field along $\gamma$ that is transverse to $P$ in the sense of Problem 10-14, then $D_t J(t) = \mathcal{H}_r(J(t))$ for all $t \in [0,b]$.
\end{exercise}

\begin{exercise}\label{prob:riccati-equation}\dcref{ch11:11-3}
Let $(M, g)$ be a Riemannian manifold, and let $f$ be any smooth local distance function defined on an open subset $U \subseteq M$. Let $F = \text{grad } f$ (so the integral curves of $F$ are unit-speed geodesics), and let $\mathcal{H}_f = \nabla F$ (the Hessian operator of $f$). Show that $\mathcal{H}_f$ satisfies the following Riccati equation along each integral curve $\gamma$ of $F$:
\[
D_t \mathcal{H}_f + \mathcal{H}_f^2 + R_{\gamma'} = 0, \quad (11.34)
\]
where $R_{\gamma'}(w) = R(w, \gamma')\gamma'$. [Hint: Let $W$ be any smooth vector field on $U$, and evaluate $\nabla_F \nabla_W F$ in two different ways.]
\end{exercise}

\begin{exercise}[Transverse Jacobi Field Comparison Theorem]\label{prob:transverse-jacobi-comparison}\dcref{ch11:11-5}
\uses{thm:jacobi-field-comparison,prob:riccati-equation,prob:jacobi-field-submanifold}
Let $P$ be an embedded hypersurface in a Riemannian manifold $(M, g)$. Suppose $\gamma: [0, b] \to M$ is a unit-speed geodesic segment with $\gamma(0) \in P$ and $\gamma'(0)$ normal to $P$, and $J$ is a normal Jacobi field along $\gamma$ that is transverse to $P$. Let $\lambda = h(J(0), J(0))$, where $h$ is the scalar second fundamental form of $P$ with respect to the normal $-\gamma'(0)$. Let $c$ be a real number, and let $u: \mathbb{R} \to \mathbb{R}$ be the unique solution to the initial value problem
\[
\begin{aligned}
u''(t) + cu(t) &= 0, \\
u(0) &= 1, \quad (11.35) \\
u'(0) &= \lambda.
\end{aligned}
\]

In the following statements, the principal curvatures of $P$ are computed with respect to the normal $-\gamma'(0)$.

\begin{enumerate}
\item[(a)] If all sectional curvatures of $M$ are bounded above by $c$, all principal curvatures of $P$ at $\gamma(0)$ are bounded below by $\lambda$, and $u(t) \neq 0$ for $t \in (0, b)$, then $|J(t)| \geq u(t)|J(0)|$ for all $t \in [0, b]$.

\item[(b)] If all sectional curvatures of $M$ are bounded below by $c$, all principal curvatures of $P$ at $\gamma(0)$ are bounded above by $\lambda$, and $J(t) \neq 0$ for $t \in (0, b)$, then $|J(t)| \leq u(t)|J(0)|$ for all $t \in [0, b]$.
\end{enumerate}

[Hint: Mimic the proof of Theorem 11.9, using the results of Problems 11-3 and 11-2.]
\end{exercise}
